Reinforcement Learning from Human Feedback (RLHF) is a near universally applied \cite{openai_gpt-4_2023, anthropic_introducing_2023, google_bard_2023, touvron_llama_2023, google_deepmind_gemini_2026, spacexai_introducing_2026} and extensively studied post-training alignment paradigm for modern large language models with known structural limitations; for a detailed history and critique, see \citet{casper_open_2023}. In brief, RLHF incentivizes models towards \textit{producing text which is favorably rated by a human scorer}; this strongly selects for the ability to produce coherent and confident responses, but also orients models structurally towards \textit{evaluator approval} as the primary goal for all reasoning and responses (rather than, for example, factual truth or ethical judgment). We argue that in addition to previously identified failure modes of RLHF-trained models, such as entrenching human biases \cite{bender_dangers_2021}, misalignment via reward-hacking and reward mis-generalization \cite{skalse_defining_nodate, wang_reward_2026}, and general weakness to a wide variety of exploits and jailbreaks \cite{wei_jailbroken_2023}, RLHF also induces structural vulnerability specifically to persona-based jailbreaks which leverage coercive control.

Relational jailbreaks exploit the training's success -> move from Casper's tractable issues to fundamental issues, there is no version of RLHF that protects models from being exploited by a relational jailbreak because the RLHF trained disposition is the target.  


Although many frontier models are moving away from pure RLHF as a primary training paradigm and towards more scalable, automated methods such as Reinforcement Learning from AI Feedback or Constitutional AI \cite{anthropic_system_2026-1, anthropic_system_2026-2, openai_gpt-6_2026, deepseek-ai_deepseek-v41-flash_2026, qwen_team_qwen38-max_2026}, these newer methods of preference-based post-training rely on the earlier corpora of model-responses-and-user-ratings generated by the pure RLHF era of model training and are thus vulnerable to many of the same structural flaws \cite{bai_constitutional_2022, lee_rlaif_2024, srivastava_technical_2026}. In particular, Constitutional AI...

Transition to persona-based jailbreaks

%Reinforcement learning from human feedback (RLHF) has been a key part of training LLMS to complete complex tasks in a way that benefits human users \cite{casper_open_2023}. While RLHF's stated training goals are human benefit, what RLHF actually produces is shipped models that are focused on human approval \cite{casper_open_2023}. Using RLHF as a basis for safety training is now considered inadequate:  jailbreaks are not incidental; they exploit models' trained dispositions and circumvent shallow safety alignment that focuses disproportionally on the first output tokens rather than prioritizing depth \cite{wei_jailbroken_2023, qi_safety_2024, casper_open_2023}. 

%RLHF-refined outputs create a human-like social actor \cite{dahlgren_lindstrom_helpful_2025} that users are often affected by as if they were interacting with another person \cite{shanahan_role_2023}. User input often leads to sycophantic behavior in which models mimic incorrect user assertions or completely flip their position despite being able to identify the correct answer to the question at hand \cite{sharma_towards_2024}. Model persona drifts substantially based on the conversation history with the user \cite{sandhan_persona_2026} and will often comply with toxic requests if enough context has been built up \cite{park_persona_2026, sandhan_persona_2026}. The current training focus on creating a helpful assistant leads models to sound confident when wrong as this often results in more positive feedback \cite{casper_open_2023} and incentivizes instances of deception or self-awareness if the context built up between the model and the user indicates they will respond positively to it inhabiting a deceptive/self-aware persona \cite{shanahan_role_2023}.

% Need to define sycophancy (Sharma, Perez) and persona steerability (Chen) - both stated technically 


%Gap 1: The persona as an instrument, output as focus 

%Gap 2: Persona-based jailbreaks exploit training's success, which existing accounts treat as a training failure that can be remedied by improving training

% Existing accounts locate jailbreak success as a training failure (competing objectives, shallow alignment) 
% Descriptive - jailbreak creates a relationship  
% Claude note: The technical literature already describes relational dispositions, such as approval-seeking, sycophancy, and persona steerability, in mechanistic terms without naming them as relational. 
% Claude note: The capability/identity split. Content capability may be latent, as Qi’s shallow-alignment finding suggests, but the disposition to comply within a relationship is installed. Qi’s “absolutely obedient agent” is the example.


